\section{Finite exclusions and a modular-jet family}
\label{sec:C-certified}
\subsection{Certified finite exclusion in the level-four family}
Put $U_C=\pi F_{1/2}(64/C)$ and
$V_C=\pi\vartheta F_{1/2}(64/C)$.
The following elementary lattice observation turns bounded numerical
searches into exact finite exclusions.

\begin{lemma}[A certificate excluding small integer relations]
Let $x_1,\ldots,x_n\in\mathbb R$, and suppose that integers
$a_1,\ldots,a_n,Q$, with $Q>0$, satisfy
$|a_i-Qx_i|\le1$. Let $\Lambda\subset\mathbb Z^{n+1}$ be the lattice
with row basis $[I_n\mid a]$. If an integer basis of $\Lambda$ has every
Gram--Schmidt squared length strictly larger than $n(n+1)H^2$, then
\[
 \sum_{i=1}^n h_ix_i=0,\qquad h_i\in\mathbb Z,
 \quad\max_i|h_i|\le H
 \quad\Longrightarrow\quad h_1=\cdots=h_n=0.
\]
\end{lemma}
\begin{proof}
A nonzero relation would yield the nonzero lattice vector
$(h_1,\ldots,h_n,\sum_i h_ia_i)$, whose last coordinate has absolute
value at most $nH$ and whose squared norm is therefore at most
$n(n+1)H^2$. For any integer basis, the squared norm of a nonzero lattice
vector is at least the least Gram--Schmidt squared length: take the
last nonzero integer basis coefficient and project onto the corresponding
orthogonalized direction. This is a contradiction.
\end{proof}

\begin{proposition}[Finite exclusion certified by exact arithmetic]
For every integer $C$ with $65\le|C|\le1024$ and
$C\notin\{-512,256\}$, no nonzero polynomial
$P\in\mathbb Z[X,Y]$ of total degree at most two and coefficient height
at most $10^{11}$ vanishes at $(U_C,V_C)$.
Consequently there is no identity
\[
 (A+B\vartheta)F_{1/2}(64/C)=\alpha/\pi
\]
in that denominator range with $(A,B)\ne(0,0)$,
$|A|,|B|\le1000$, and $\alpha$ annihilated by a nonzero integer
polynomial of degree at most two and coefficient height at most $50000$.
\end{proposition}
\begin{proof}[Proof with a reproducible finite certificate]
For each of the $1918$ indicated parameters, enclose the six monomials
\[
 U_C^2,\ U_CV_C,\ V_C^2,\ U_C,\ V_C,\ 1
\]
using outward-rounded integer intervals of scale $10^{120}$. These
enclosures follow from
\[
 \frac{\binom{2n+2}{n+1}^3|C|^{-n-1}}
      {\binom{2n}{n}^3|C|^{-n}}
 =\frac{8(2n+1)^3}{(n+1)^3|C|}\le\frac{64}{|C|}
\]
and geometric bounds for the tails of both the series and its Euler
derivative. An exact alternating-series enclosure from Machin's formula
encloses $\pi$. With $Q=10^{80}$, nearest-midpoint rounding gives
integers $a_i$; both endpoint inequalities $|a_i-Qx_i|\le1$ are checked
exactly. Every stored integral basis is obtained by a matrix of determinant
$\pm1$, with the matrix identity checked. Independent rational
Gram--Schmidt computations show in all $1918$ cases that the least squared
length exceeds $42(10^{11})^2$.

The accompanying scripts and the parameter-by-parameter integer bases are
stored in the research archive under the stem
\texttt{20260918C\_numeric\_level4\_lattice}; verifying a stored basis
does not require trusting lattice reduction. The lemma proves the first
assertion. Substituting $AX+BY$ into an annihilating polynomial for
$\alpha$ gives a nonzero polynomial in $X,Y$, of total degree at most two
and coefficient height at most
$2\cdot1000^2\cdot50000=10^{11}$, proving the second.
\end{proof}

This is a bounded arithmetic exclusion, not a proof of CM forcing with
unrestricted coefficients or algebraic multipliers. The three analogous
tests at $C=-512,256,4096$ correctly fail to exclude their known small
polynomial relations.

\subsection{Additional exact finite searches}
The interval method also gives exclusions for multiplier sets with
unbounded numerator.  At level four, take
\[
 65\leq |C|\leq10000,\qquad |A|\leq20,\quad 0\leq B\leq100,
 \qquad \gcd(A,B)=1,
\]
using $(A,B)=(1,0)$ as the representative when $B=0$.
There are $19872$ denominators and $2500$ coefficient pairs, hence
$49680000$ cases.  If $\alpha$ is rational of reduced denominator at
most $100$, the only surviving identities are the known rows
$(1,6,256)$ and $(5,42,4096)$.  If instead $\alpha^2$ is rational
of reduced denominator at most $100$, the only additional row is
$(1,6,-512)$.  Each surviving row is already established exactly.

For each of levels one to three, take the $127$ denominators
\[
 \{R_\ell+1,\ldots,R_\ell+63\}
 \ \cup\ \{-R_\ell,-R_\ell-1,\ldots,-R_\ell-63\},
\]
and the $502$ primitive representatives with $|A|\leq10$,
$1\leq B\leq40$, together with $(1,0)$.  No identity in this
set has either $\alpha^2$ or $\alpha^4$ rational with reduced
denominator at most $1000$.  These are $382524$ power tests.
The exact CM class-polynomial certificate verifies that every
parameter in this last set is non-CM.

Here are the certification principles. Rational recurrences and
geometric tail bounds give outward enclosures for the series and
its Euler derivative. Near the boundary the quadratic Clausen
parameter has absolute value at most $1/2$, so its Gauss series
gives rapidly convergent enclosures, including the negative
endpoint. Machin's alternating arctangent formula supplies an
independent rational enclosure for $\pi$. For an enclosure of
$\alpha^k$ and a denominator bound $M$, an exact rational
nearest-fraction calculation decides whether the enclosure
contains a reduced fraction with denominator at most $M$.
The enclosures are narrower than $1/M^2$, the minimum separation
between two such distinct fractions. Exclusion therefore does
not require an upper bound on their numerators.

All endpoints and comparisons are rational. A separate integer
implementation recomputes the enclosures and verifies the rational
detector. The scripts and certificates, including their input and
source hashes, are stored under
\texttt{20260918C\_numeric\_level4} and
\texttt{20260918C\_numeric\_levels123} in the attempt archive.

Higher-degree numerical integer-relation searches likewise found
only known CM controls. Their null results are evidence only and
are not used in any theorem or certificate above.

\subsection{An explicit non-CM test family and an exact infinite exclusion}
\label{subsec:C-explicit-jet}

The affine two-jet formulation defines a family in which a counterexample
would have to meet an additional arithmetic condition. The rational
constants of bounded denominator below can be excluded unconditionally.
Put $J(y)=j(iy)$ for $y>1$, with primes on $J$ denoting real $y$-derivatives.

\begin{proposition}\label{prop:C-jet-family}
The function $H(y)=J''(y)/J'(y)$ is a strictly decreasing bijection from
$(1,\infty)$ to $(2\pi,\infty)$. For $c>2\pi$, let $H(y_c)=c$,
$C_c=J(y_c)$, and $t_c=\sqrt{1-1728/C_c}$. Then $C_c>1728$ is strictly
decreasing in $c$, and
\begin{equation}\label{eq:C-jet-family-identity}
 \bigl[(864-C_c)+(C_c-1728)\thetaop\bigr]F_{1/6}(1728/C_c)
 =-\frac{c(C_c-1728)}{2\pi t_c}.
\end{equation}
The series converges absolutely. If $c$ is algebraic, its elliptic fiber is
non-CM. Nevertheless, $C_c\notin\mathbb Z$ for every integer $c>2\pi$.
\end{proposition}

\begin{proof}
Writing $q=e^{-2\pi y}$, positivity of the coefficients in
$qj=E_4^3\prod_{n\ge1}(1-q^n)^{-24}$ gives
\[
 J=q^{-1}+744+\sum_{n\ge1}a_nq^n,\quad a_n>0,\qquad
 U=\sum_{n\ge1}na_nq^{n+1},\quad
 V=\sum_{n\ge1}n^2a_nq^{n+1}.
\]
Consequently $J'=2\pi q^{-1}(1-U)$ and
$H=2\pi(1+V)/(1-U)$. Since $j'(i)=0$, one has $U(1)=1$.
Both $U$ and $V$ strictly decrease with $y$, proving that $H$ strictly
decreases from infinity to $2\pi$. Also $J$ strictly increases from $1728$
to infinity. These facts prove existence, uniqueness, and monotonicity.

For the standard level-one model $g_2=3$, $g_3=t$,
\[
 F_{1/6}=\frac{3\omega^2}{2\pi^2},\qquad
 \frac{\thetaop F_{1/6}}{F_{1/6}}
 =-\frac16-\frac{h}{3t\omega}.
\]
Thus the quadratic period coefficients for $(A+B\thetaop)F_{1/6}$ are
$a=3A/2-B/4$ and $b=-B/(2t)$. The period--jet dictionary gives
\[
 bj''+\left(\frac{a}{3jt}-bR(j)\right)j'^2+i\alpha j'=0,
 \qquad R(j)=\frac2{3j}+\frac1{2(j-1728)}.
\]
Its $j'^2$ coefficient vanishes precisely when
$A/B=-(j-864)/(j-1728)$. Since $j''(iy_c)/j'(iy_c)=-ic$,
choosing $A=864-C_c$, $B=C_c-1728$ yields
$\alpha=bc=-c(C_c-1728)/(2t_c)$, proving
\eqref{eq:C-jet-family-identity}.

To exclude CM when $c$ is algebraic, use the exact formula
\[
 \frac{j''}{j'}=
 \left(R(j)-\frac{s_2}{6j}\right)j'+\frac{i}{y},\qquad
 s_2=\frac{E_4}{E_6}\left(E_2-\frac3{\pi y}\right).
\]
At a CM point on this branch, $j,s_2,y$ are algebraic, while $j'$ is
transcendental. Here $E_6\ne0$, so the algebraicity assertion for $s_2$
in \cite[Theorem~8.1]{CGG} applies. Comparing with $j''/j'=-ic$ would
force $-c=1/y$, which is impossible.

It remains to exclude integer $C_c$ for integer $c$. This last claim uses
a finite exact certificate. The archived verifier
\texttt{20260918C\_explicit\_jet\_search.py} and its accompanying JSON file
contain rational brackets $q_-<q_+$ for all $c=7,\ldots,200$.
The brackets have denominator $10^{65}$. Their validity is checked by
384-bit directed dyadic interval arithmetic using
\[
 H=\frac\pi3\left(4\frac{E_6}{E_4}
       +3\frac{E_4^2}{E_6}-E_2\right),\qquad
 J=\frac{1728E_4^3}{E_4^3-E_6^2}.
\]
The Eisenstein sums are truncated after $50$ terms. For $k=1,3,5$ and
$0<q<1/500$, their omitted sums satisfy the elementary bound
\[
 0\le\sum_{n>50}\sigma_k(n)q^n
 \le2\cdot51^{k+1}q^{51},
\]
since $\sigma_k(n)\le n^{k+1}$ and the successive majorant ratios are
at most $(1/500)(52/51)^{k+1}<1/2$. A rational enclosure of $\pi$ follows
from Machin's formula using $110$ alternating terms in each arctangent.
Every bracket has certified $E_6>0$, so it lies on the required branch.
For each of the $194$ brackets, the verifier establishes
$H(q_-)<c<H(q_+)$ and places
$[J(q_+),J(q_-)]$ strictly between consecutive integers. Approximate
arithmetic proposes the brackets only; their replay uses exact integers.
The smallest certified distance to an integer over this range exceeds
$0.0032177275685462482644977288$, at $c=156$.
In particular,
\[
 1728.6093005038821554072482832073<C_{200}
 <1728.6093005038821554072482832074<1729.
\]
Finally, every integer $c>2\pi$ is at least $7$, and monotonicity gives
$1728<C_c\le C_{200}<1729$ for $c\ge200$. The finite certificate and
this last inequality cover all integer $c$.
\end{proof}

The proposition excludes a genuine candidate mechanism, rather than a
normalization defect. If an algebraic noninteger $c>2\pi$ had integer
$C_c$, division of $A=864-C_c$ and $B=C_c-1728$ by
$\gcd(C_c-864,C_c-1728)$ in \eqref{eq:C-jet-family-identity} would yield
a primitive standard identity at a non-CM fiber. Algebraicity of such a
$c$ remains unexcluded. Equivalently, the unconditional result above
proves $H(y(C))\notin\mathbb Z$ for integer $C>1728$; it does not prove
$H(y(C))\notin\overline{\mathbf Q}$.

\begin{proposition}\label{prop:C-rational-jet}
In the notation of Proposition~\ref{prop:C-jet-family}, if $C_c$ is an
integer and $c=a/b$ is rational in lowest terms, with $b>0$, then $b>5$.
More generally, for all rational $c=a/b$ in this family,
\[
 b\gg C_c^{125/444}\qquad(C_c\longrightarrow\infty).
\]
Consequently, for fixed $L>0$ and $\gamma<125/444$, only finitely many
integer fibers in this family can have $b\le L C_c^\gamma$.
The implied constant is not made effective here.
\end{proposition}

\begin{proof}
The certificate \texttt{20260918C\_jet\_rational.json} contains all
$1938$ reduced rational numbers $c$ with $1\le b\le5$ and
$2\pi<c\le200$. The endpoint list is determined using the exact Machin
bounds from Proposition~\ref{prop:C-jet-family}; its smallest element
is $19/3$. The accompanying verifier checks every rational root bracket
by the dyadic method of that proposition and by a separately implemented
decimal-grid interval method. The latter uses $45$ direct Eisenstein
terms and the uniform tail $2\cdot46^{k+1}/500^{46}$ for $k=1,3,5$.
Both verifiers prove that each $C_c$ lies strictly between consecutive
integers. This adds $1744$ nonintegral rational constants to the earlier
$194$ integer constants. The smallest certified distance to an integer
is greater than
\[
 0.000048699599821220308647971801,
\]
at $c=623/4$. In particular,
\[
 1728.9999513004001787796913520281<C_{623/4}
 <1728.9999513004001787796913520282<1729.
\]
Monotonicity now excludes integer $C_c$ for every real $c\ge623/4$;
together with the finite certificate, this proves $b>5$.

For the growth statement, the cusp expansion
\[
 J=q^{-1}+744+196884q+O(q^2)
\]
and the functions $U,V$ in Proposition~\ref{prop:C-jet-family} give
$U=196884q^2+O(q^3)$ and $V=196884q^2+O(q^3)$. Therefore
\[
 H=2\pi\frac{1+V}{1-U}
   =2\pi+787536\pi q^2+O(q^3).
\]
Since $q=C_c^{-1}+O(C_c^{-2})$, a rational value $H=a/b$ satisfies
\[
 0<\frac{a}{2b}-\pi
   =\frac{393768\pi}{C_c^2}+O(C_c^{-3}).
\]
The denominator of $a/(2b)$ in lowest terms is at most $2b$.
Applying the Zeilberger--Zudilin bound
\eqref{eq:pi-rational-lower-bound}, with $\mu_*=7.104=888/125$, yields
\[
 c_*(2b)^{-888/125}\ll C_c^{-2},
\]
and hence $b\gg C_c^{125/444}$. If
$b\le L C_c^\gamma$ with $\gamma<125/444$, this bounds $C_c$;
integrality and the uniqueness of $c$ at each fiber prove finiteness.
\end{proof}

The primitive slopes in this family have a simple description. Put
$d=\gcd(C_c-1728,864)$ and $k=864/d$. Then
\[
 (A,B)=(-B-k,B),\qquad
 C_c=1728+\frac{864}{k}B,\qquad \gcd(B,k)=1,
\]
where $k$ ranges over the $24$ positive divisors of $864$.
For rational $c$, the proposed multiplier satisfies
$\alpha^2=c^2C_c(C_c-1728)/(4d^2)\in\Q$.
Thus this restriction addresses a specific part of the quadratic
multiplier case. It does not exclude rational $c$ of arbitrary
denominator, or algebraic irrational $c$.

